\chapter{Assembly and verification}
\label{ch:assembly}

\section{The full operator}

The full model maps four unknowns to two observables:
\[
    (\delta\rho_m,\ x_L,\ h_s,\ \delta K)\ \longmapsto\ (h,\ N) .
\]
For each source $X$ of load (mantle density, lithospheric unknowns, the sea-level constant), with load operator $L_X$ and direct geoid $D_X$,
\begin{align}
    N_X &= (I+T)^{-1}\left(D_X + G_w P L_X\right) && \text{the geoid, ocean included},\\
    w_X &= P\left(L_X - \rho_w g\,C N_X\right) && \text{the deflection under its load and the water's},\\
    h_X &= w_X - N_X && \text{topography from sea level},
\end{align}
and $h = \sum_X h_X x_X + h_s + h_{\mathrm{ref}}$, $N = \sum_X N_X x_X + N_{\mathrm{ref}}$.
The reference terms $(h_{\mathrm{ref}}, N_{\mathrm{ref}})$ are the same equations applied to WINTERC-G's own load and structure: they are the fixed offset of the affine forward model.
The surface field $h_s$ is added to the topography with no load and no geoid: it is a nuisance field for short-wavelength topography that the long-wavelength model is not meant to explain, separated from the rest only by its prior, which puts its power at high degree.

\begin{incode}{\texttt{src/density\_inversion/forward/model.py}, \texttt{ForwardModel}}
\begin{itemize}[nosep]
\item \texttt{model\_space}, \texttt{data\_space}: direct sums of pygeoinf spaces, labelled.
\item \texttt{forward(x)}: $(h, N)$ including the offsets; \texttt{offset()}: the offsets alone.
\item \texttt{parts(x)}: every contribution separately (topography by source, geoid by mass), which sum to \texttt{forward}.
\item Each block is a pygeoinf \texttt{LinearOperator} with an adjoint, so the posterior can be computed matrix-free; the adjoints are consistent to $1.3\times10^{-11}$.
\end{itemize}
\end{incode}

\section{The verification}

Each check compares the code with something it does not share code with.

\begin{center}
\begin{longtable}{p{0.33\textwidth}p{0.42\textwidth}r}
\toprule
Check & Against & Agreement\\
\midrule
\endhead
Sheet factor $S_\ell$, force balances & symbolic derivation (sympy) & $3\times10^{-15}$\\
Core model & closed form & $4\times10^{-14}$\\
Topography kernels, $\ell=1$--40 & \texttt{fem\_1d}, separate finite-element code & \qty{0.09}{\percent}\\
Geoid kernels, $\ell=1$--40 & \texttt{fem\_1d} & \qty{0.1}{\percent}\\
Degree 1 & \texttt{fem\_1d} at $\ell=1$ & $10^{-4}$\\
Full model in the core limit & flowkernels' $K_N^{\mathrm{fk}}$, $K_{h0}$ & $10^{-10}$\\
Plate & manufactured solution & $10^{-9}$\\
Lithospheric load & Airy isostasy & \qty{0.06}{\percent}\\
Variable-depth sheets & brute-force integration & $10^{-9}$\\
Moved mass & mass balance (load = geoid sheets) & $10^{-9}$\\
Reference column & exact column mass & $10^{-12}$\\
Grid quadrature & refinement to $8L$ & \qty{0.2}{\percent}\\
Adjoints & $\langle Ax, y\rangle = \langle x, A^*y\rangle$ & $1.3\times10^{-11}$\\
Compensated root geoid & brute-force axisymmetric code & \qty{3}{\percent}\\
\bottomrule
\end{longtable}
\end{center}

The last row is the weakest.
The brute-force code solves the Airy balance column by column and integrates Newton's law numerically, sharing nothing with the package.
The cause of the \qty{3}{\percent} gap ($-0.240$ against $-0.248$\,m) has not been established.
Planar versus spherical mass balance is one candidate, but in the audit case it changed the result by about \qty{15}{\percent} ($-0.254$ against $-0.219$\,m), which is not obviously consistent with a \qty{3}{\percent} gap.
Until the gap is shown to shrink with resolution, describe it as an unexplained \qty{3}{\percent} disagreement, not rounding error.

\section{The approximations}

\begin{center}
\begin{longtable}{p{0.4\textwidth}p{0.45\textwidth}}
\toprule
Approximation & Size\\
\midrule
\endhead
Covariant flat plate vs thin shell & $<10^{-4}$ in $F_\ell$ for $\ell\le50$\\
Membrane stresses neglected & \qtyrange{0.2}{0.4}{\percent}\\
One $g$ in the plate and loads & \qty{0.4}{\percent} at sea\\
Radial grid of the flow kernels & $6\times10^{-4}$; up to \qty{9}{\percent} for a deep anomaly at $\ell=40$\\
First-order lithosphere & \qty{4}{\percent} at $\ell=50$ for \qty{10}{km}; cross terms \qtyrange{3}{11}{\percent}\\
Lid base fixed at \qty{105}{km} & cratons \qty{+865}{m} too high in the reference\\
One crustal layer & up to \qty{1.3}{m} of geoid\\
Ice density 910 vs 950 & documented\\
Phase-boundary deflections (410, 660) & not modelled; raised by a referee report\\
\bottomrule
\end{longtable}
\end{center}

\begin{defence}
\Q{Your radial grid converges only to first order. Is that a problem?}
\A For the topography and for degrees up to about 10 the kernels are converged to $6\times10^{-4}$.
For deep, sharp anomalies at high degree the error reaches \qty{9}{\percent}, but those combinations produce almost no surface signal.
Refinement or Richardson extrapolation is the remedy if they become important.

\Q{What about the 410 and 660 phase transitions?}
\A Their deflections are mass anomalies that the model does not represent explicitly.
To first order they behave like additional density sheets at those depths, which the mantle unknown can partly absorb; whether the prior allows it is something to test with synthetics.
\end{defence}
